\par\bigskip\begingroup\small
\begin{center}\textls[120]{Erläuterungen zum Briefwechsel}, S. 429—437.\\
\textls[120]{Geschichtliches zu den Abschnitten} [V.], [VI.], [VII.].\end{center}
\indent In den Briefstellen [3.] und [4.] schreibt \textsc{Gauss}, daß er in den ersten Jahren des XIX. Jahrhunderts, beziehungsweise etwas mehr oder weniger als 50 Jahre vor dem 5. Februar 1850 eine Methode gefunden habe, um den »Grad der Konvergenz« trigonometrischer Reihen, insbesondere der Entwicklung der Mittelpunktsgleichung, zu bestimmen. Auch habe er damals die Richtigkeit dieser Konvergenzuntersuchung durch numerische Rechnung an einem Beispiel der Mittelpunktsgleichung für einen großen Wert der Exzentrizität bestätigt; seine »Konvergenzformel« für die Koeffizienten der Entwicklung der Mittelpunktsgleichung stimme mit der von \textsc{Jacobi} (siehe das Zitat oben S. 432 Fußnote **)) gegebenen überein.

\indent Dazu ist zuvörderst zu bemerken, daß der Ausdruck »Grad der Konvergenz« hier im Sinne des art. 1. von Abschnitt [V.], oben S. 400, den Grad der Annäherung der Reihenkoeffizienten an die Null bedeutet, wie denn auch die in der Briefstelle [3.] aus der »Konvergenzformel« gezogene Folgerung, nämlich die Vergleichung der Konvergenz der Reihenkoeffizienten mit der einer geometrischen Progression, ganz im Sinne dieses art. [1.] gehalten ist.

\indent Die Feststellung vom 5. Februar 1850 macht \textsc{Gauss} auf Grund einer von ihm damals aufgefundenen Aufzeichnung, die jene numerische Rechnung enthielt. Nun finden wir auf dem zweiten Zettel, der in den artt. [2.] und [3.] des Abschnitts [VII.] abgedruckt ist, einmal die mit der \textsc{Jacobi}schen völlig übereinstimmende
\clearpage\SourceLeaf{450}\ArithmeticHead{443}{R}{BEMERKUNGEN ZUR REIHENLEHRE. BRIEFWECHSEL.}
\noindent asymptotische Formel für die Entwicklungskoeffizienten der Mittelpunktsgleichung\textsuperscript{*}); ferner ist die Tabelle des art. [3.], wie wir oben S. 441 gesehen haben, für den Wert \(0{,}6550518\) der Exzentrizität gerechnet, und dieser Wert ist in der Tat als \textls[120]{groß} zu bezeichnen (siehe oben S. 433, erste Zeile), da er die Exzentrizitäten der großen Planeten erheblich übertrifft und schon nahe an der \textsc{Laplace}schen Grenze \(0{,}6627430\,{.}\;{.}\)\textsuperscript{**}) liegt. Auch die Angabe in der Briefstelle [3.], daß die Herleitung der Formel eine halbe Oktavseite einnehme, trifft für unsern Zettel wenn auch nicht buchstäblich, so doch dem Sinne nach zu, da die Entwicklung, die zu dem Ergebnis führt, in der Tat sehr kurz ist. Die Bemerkung der Briefstelle [4.], daß die Methode auch geeignet sei, alle Glieder numerisch zu berechnen, nicht nur die entfernten, und zwar mit jeder Genauigkeit, findet auf dem dritten Zettel ihre Bestätigung (siehe unsern art. [4.]); die auf diesem Zettel befindlichen nicht mit abgedruckten Zahlenrechnungen beziehen sich auch auf das im art. [3.] behandelte Beispiel, sind aber nicht zu Ende geführt, was mit der Briefstelle [4.] im Einklang steht.

\indent Es ist also unzweifelhaft, daß wir in den beiden Zetteln, die in den artt. [2.], [3.], [4.] abgedruckt sind, die von \textsc{Gauss} 1850 wiedergefundenen Aufzeichnungen vor uns haben. Was ihre Entstehungszeit anlangt, so können wir das »\(50\pm\)« des Briefes [4.] in Übereinstimmung mit der Angabe des Briefes [3.] genauer zu »\(50-\)« bestimmen, d. h. also die Abfassung \textls[120]{dieser} Aufzeichnungen in die ersten Jahre des XIX. Jahrhunderts setzen. Denn erstens bedient sich \textsc{Gauss} darin stets der Bezeichnung \(i\) für \(\sqrt{-1}\), was er im XVIII. Jahrhundert noch nicht getan hat, zweitens tritt am Schluß des art. [3.] die Darstellung des konstanten Teils von \((1-g\cos M)^{-\frac12}\) durch das arithmetisch-geometrische Mittel auf, und diese Darstellung hat \textsc{Gauss} erst um den 23. Dezember 1799 gefunden (siehe oben S. 183, 185 und die zugehörigen Bemerkungen S. 273—274). In dem Beispiel des \textls[120]{ersten} Zettels, oben S. 421, ist der Wert der Exzentrizität \(\varepsilon=0{,}254236\) den V. Elementen des kleinen Planeten Juno entnommen, die in der Monatlichen Correspondenz \textbf{11}, S. 476, Mai 1805 (Werke VI, S. 264) veröffentlicht sind, und die \textsc{Gauss} gestützt auf seine Beobachtung vom 20. Februar 1805 berechnet hatte. Die VI. Elemente der Juno, die \textsc{Gauss} am 25. August 1806 an die Monatliche Correspondenz gesandt hat (siehe Monatl. Corr. \textbf{14}, S. 378, Werke VI, S. 280, vergl. auch S. 279), geben die Exzentrizität \(0{,}254944\). Der erste Zettel ist also jedenfalls zwischen dem 20. Februar 1805 und dem 25. August 1806 geschrieben. Wahrscheinlich stammt auch der \textls[120]{vierte} Zettel aus dieser Zeit; \textsc{Gauss} würde also den Kunstausdruck »Zug komplexer Werte«, den er im Abschnitt [VI.] art. 2., oben S. 408, erklärt, schon zu Anfang des XIX. Jahrhunderts angewandt haben, was immerhin Rückschlüsse auf seine damalige Bewandertheit in der Analysis der komplexen Veränderlichen gestattet.

\indent Es sind aber noch andere, bis ins XVIII. Jahrhundert zurückreichende Spuren vorhanden, die auf \textsc{Gauss}ens Beschäftigung mit diesem Gegenstande hinweisen und die auch darüber hinaus einen Einblick in den ganzen Gedankenkreis gewähren, der für \textsc{Gauss} in diesem Zusammenhang in Betracht kam.

\indent Den Ausgangspunkt bildete für \textsc{Gauss} die noch in die Göttinger Studentenjahre fallende Beschäftigung mit dem \textsc{Lagrange}schen Lehrsatz, die in der Aufzeichnung Nr. 49 des \textit{Tagebuchs} vom 27. Dezember 1796 angezeigt ist, und für die der Brief [1.] an \textsc{Hindenburg}\textsuperscript{***}), sowie die Werke VIII, S. 76 abgedruckte
\par\medskip\noindent\rule{.24\linewidth}{.4pt}\par\smallskip
\noindent\hspace*{\parindent}\textsuperscript{*}) Um den \textsc{Jacobi}schen Ausdruck (siehe \textsc{Jacobis} Werke VII, S. 188 und 237)
\[
\left\{\frac{e i^{\sqrt{1-ee}}}{1+\sqrt{1-ee}}\right\}^{p}
\left(\frac1p+\frac4{3\sqrt{2p^3\pi}}\frac1{\sqrt[4]{(1-ee)^3}}\right)
\]
\noindent mit dem \textsc{Gauss}schen in Übereinstimmung zu bringen, hat man zu setzen \(f=\sin\varphi\) für \(e\), \(e\) für \(i\) und \(n\) für \(p\); dann ergibt sich unmittelbar unsere Formel [28]$'$, oben S. 441.
\par\noindent\hspace*{\parindent}\textsuperscript{**}) \textsc{Laplace}, \textit{Mécanique céleste}, t. V, 1825, Suppl. \textit{Sur le développement des coord. ellipt.} S. 11.
\par\noindent\hspace*{\parindent}\textsuperscript{***}) Die Urschrift dieses Briefes befindet sich im Privatbesitz; dem Abdruck liegt eine im \textsc{Gauss}archiv befindliche, von C. H. \textsc{Müller} angefertigte Abschrift zugrunde.
\par\smallskip\hfill 56\textsuperscript{*}
\clearpage\SourceLeaf{451}\ArithmeticHead{444}{L}{ANALYSIS. NACHLASS.}
\noindent Abhandlung mit der zugehörigen Bemerkung von R. \textsc{Fricke} zu vergleichen ist\textsuperscript{*}). \textsc{Gauss} wird wohl die Abhandlung von \textsc{Lagrange} \textit{Sur le problème de Képler}, die im 25. Bande (1769), 1771, S. 204 der Memoiren der Berliner Akademie\textsuperscript{**}) steht, auch schon um diese Zeit gekannt haben, da er die auf die Umkehrungsformel bezügliche, die im 24. Bande derselben Sammlung steht, in der Abhandlung Werke VIII, S. 76 selbst zitiert\textsuperscript{***}). In jener Abhandlung \textit{Sur le problème de Képler} gewinnt \textsc{Lagrange} die Entwicklungen der drei Größen \(E,r,v\) (in der Bezeichnung von \textsc{Gauss}) nach den Sinus bezw. Kosinus der Vielfachen der mittleren Anomalie \(M\), indem er aus den Gleichungen (I), (II), oben S. 438, und aus der Polargleichung der Ellipse durch Anwendung seiner Umkehrungsformel zuerst nach Potenzen der Größe \(f\sin M\) entwickelt, dann die Potenzen von \(\sin M\) durch die trigonometrischen Funktionen der Vielfachen von \(M\) ausdrückt und nach diesen ordnet; er setzt:
\[
\begin{aligned}
E&=M+\sum A_n\sin nM,\\
r&=1+\frac{f^2}{2}-\sum B_n\cos nM,\\
v&=M+\sum K_n\sin nM
\end{aligned}
\]
\noindent und drückt die \(A_n,B_n\) durch einfache, die \(K_n\) durch kompliziertere Reihen aus, die nach Potenzen der Exzentrizität \(f\) fortschreiten. Die Bestimmung des infinitären Verhaltens der Koeffizienten \(A_n,B_n\) kommt auf die Untersuchung der \textsc{Bessel}schen Funktion \(J_n(nf)\) und ihrer Derivierten für große Werte von \(n\) hinaus. Die oben S. 385 abgedruckte Notiz Nr. 82 des \textit{Tagebuchs} vom 16. Oktober 1797 zeigt, wie wir S. 389 nachgewiesen haben, daß \textsc{Gauss} sich etwa 10 Monate nach den Studien zur \textsc{Lagrange}schen Reihe mit der asymptotischen Darstellung der \textsc{Bessel}schen Funktion \(J_1(x)\) beschäftigt hat.
\par\medskip\noindent\rule{.24\linewidth}{.4pt}\par\smallskip
\noindent\hspace*{\parindent}\textsuperscript{*}) An der von \textsc{Fricke} zitierten Stelle: \textsc{Sartorius v. Waltershausen}, \textit{Gauss zum Gedächtniss}, 1856, S. 22, wird gesagt, \textsc{Pfaff} habe damals (also im April 1797) den Aufsatz von \textsc{Gauss} an \textsc{Hindenburg} befördert; wie der Brief [1.] zeigt, war es aber nicht \textsc{Pfaff}, sondern \textsc{Kaestner}. Auch die Angabe \textsc{Sartorius v. Waltershausens}, \textsc{Hindenburg} sei bald nach Empfang des \textsc{Gauss}schen Manuskripts gestorben, trifft nicht zu, denn \textsc{C. Fr. Hindenburg} starb in Leipzig am 17. März 1808.
\par\noindent\hspace*{\parindent}\textsuperscript{**}) Oeuvres de \textsc{Lagrange} III, S. 113.
\par\noindent\hspace*{\parindent}\textsuperscript{***}) Unter den im Nachlaß befindlichen \textls[120]{Auszügen} aus Zeitschriften u.s.w. (Ca4) sind Bemerkungen zu den in den Bänden 1—25 der Berliner \textls[120]{Mémoires} enthaltenen mathematischen Abhandlungen vorhanden.
\par\endgroup
\ArithmeticHead{444}{L}{ANALYSIS. NACHLASS.}
